\documentclass[main.tex]{subfiles}

\begin{document}

\subsection{Constructions of Concrete Tensor Product Structures}
\label{subsec:tensor_products_of_operator_systems_II}

% Organize the equality problems among the minimal, maximal, commuting, and
% injective tensor products.  State the corresponding nuclearity properties
% and isolate the comparisons relevant to tensor-product and commuting-
% operator models in self-testing.

With the abstract work of the previous subsection squared away, we will now move on to giving explicit constructions of tensor product structures that will be of interest to us. In Lecture 2, Hridoyananda provided the minimal and maximal tensor product structures of 
Constructions \ref{const: minimal tensor product structure} and \ref{const: maximal tensor product structure}.

\begin{construction}{Minimal Tensor Product Structure}{const: minimal tensor product structure}

\end{construction}

\begin{construction}{Maximal Tensor Product Structure}{const: maximal tensor product structure}

\end{construction}

The next two constructions, the commuting \& joint tensor product structures, were provided by Lucas in Lecture 3. 

\begin{construction}{Commuting Tensor Product Structure}{const: commuting tensor product structure}

\end{construction}

\begin{construction}{Joint Tensor Product Structure}{const: joint tensor product structure}

\end{construction}


Of course, in these constructions, we have provided \emph{some} structure that we have referred to as an operator system tensor product but the next step is to prove that these actually do give rise to tensor product structures. The proofs are actually rather technical 


\end{document}
